% !TEX root = task-2-5-1.tex
% Editable task bank. Append new BankTask blocks at the end of the document.
\documentclass[12pt,a4paper]{article}
\usepackage{physics-textbook}
\usepackage{graphicx,needspace}
\usetikzlibrary{decorations.pathmorphing}
\newif\ifShowSolutions
\ifdefined\TutorEdition
  \ShowSolutionstrue
\else
  \ShowSolutionsfalse % Change to \ShowSolutionstrue to display solutions.
\fi
\newcounter{banktask}
\newsavebox{\BankTaskBox}
\newcommand{\BankTaskContent}[3]{%
  \noindent\textbf{\thebanktask.}\quad #2\par
  \ifShowSolutions
    \Needspace{50mm}% Keep the solution heading with its opening text or diagram.
    #3\par
    \ifdefined\TutorEdition
      \ifcsname TutorSource#1\endcsname
        \nopagebreak\smallskip{\small\csname TutorSource#1\endcsname\par}%
      \fi
    \fi
  \fi
}
\newcommand{\BankTask}[3]{%
  \refstepcounter{banktask}\par\medskip
  \sbox{\BankTaskBox}{\begin{minipage}{\linewidth}%
    \BankTaskContent{#1}{#2}{#3}\end{minipage}}%
  \ifdim\dimexpr\ht\BankTaskBox+\dp\BankTaskBox\relax>\dimexpr\textheight-20mm\relax
    \Needspace{35mm}\BankTaskContent{#1}{#2}{#3}%
  \else
    \Needspace{\dimexpr\ht\BankTaskBox+\dp\BankTaskBox\relax}%
    \noindent\usebox{\BankTaskBox}\par
  \fi
  \medskip
}
\begin{document}
\BookSubsection{2.5.1}{Средняя кинетическая энергия частиц и температура газа}


% Source sdamgia-817: https://phys-ege.sdamgia.ru/problem?id=817
\BankTask{sdamgia-817}{
В результате нагревания идеального газа средняя кинетическая энергия теплового движения его молекул увеличилась в $4$ раза. Во сколько раз изменилась абсолютная температура газа?
}{
\smallskip\noindent\textbf{Решение}\par
Средняя кинетическая энергия теплового движения молекул идеального газа прямо пропорциональна абсолютной температуре: 

$\overline{E}= \frac{3}{2} kT.$

Следовательно, при увеличении средней кинетической энергии теплового движения в $4$ абсолютная температура газа также увеличится в $4$ раза.

Ответ: $4$.
}

% Source sdamgia-9735: https://phys-ege.sdamgia.ru/problem?id=9735
\BankTask{sdamgia-9735}{
При построении температурной шкалы Реомюра принимается, что при нормальном атмосферном давлении лед тает при температуре $0$ градусов Реомюра (°$R$), а вода кипит при температуре $80^\circ$R. Найдите, чему равна средняя кинетическая энергия поступательного теплового движения частицы идеального газа при температуре $91^\circ$R. (В электрон-вольтах; округлите до сотых долей.)
}{
\smallskip\noindent\textbf{Решение}\par
Из описания следует, что один градус шкалы Реомюра соответствует $\frac{100}{80} =1{,}25 ^\circ\text{C} .$

Соответственно, $91 ^\circ R$ — это $100 + 11 \cdot 1{,}25 = 113{,}75 ^\circ\text{C}$ или $113{,}75 + 273=386{,}75\text{К}.$

Средняя кинетическая энергия теплового движения частиц идеального газа связана с температурой соотношением: 

$\overline{E}_{\text{к}} = \frac{3}{2} kT=1{,}5 \cdot 1{,}38 \cdot 10 ^{- 23} \text{Дж}/\text{К} \cdot 386{,}75\text{К}\approx8 \cdot 10 ^{- 21} \text{Дж}\approx 0{,}05\text{эВ}.$ 

Ответ: $0{,}05$.
}

% Source sdamgia-11663: https://phys-ege.sdamgia.ru/problem?id=11663
\BankTask{sdamgia-11663}{
Кислород и водород находятся в закрытом сосуде в состоянии термодинамического равновесия друг с другом. Во сколько раз среднеквадратичная скорость молекул водорода отличается от среднеквадратичной скорости молекул кислорода?
}{
\smallskip\noindent\textbf{Решение}\par
Средняя энергия теплового движения молекул связана с абсолютной температурой газа соотношением: 

$\overline{E}= \frac{m\overline v ^{2}}{2} = \frac{3}{2} kT,$ 

где $m$ — масса одной молекулы.

По определению среднеквадратичная скорость равна: 

$\bar v = \sqrt{\overline v ^{2}} = \sqrt{\frac{3RT}{M}} .$ 

Значит, отношение среднеквадратичных скоростей молекул водорода и кислорода:

$\frac{\bar{v}_{\mathrm{H}_2}}{\bar{v}_{\mathrm{O}_2}}=\sqrt{\frac{M_{\mathrm{O}_2}}{M_{\mathrm{H}_2}}}=\sqrt{\frac{32}{2}}=4.$

Ответ: $4$.
}

% Source sdamgia-12961: https://phys-ege.sdamgia.ru/problem?id=12961
\BankTask{sdamgia-12961}{
При некотором значении среднеквадратичной скорости хаотического движения молекул идеального газа средняя кинетическая энергия его молекул равна $56\cdot 10^{-22}\,\text{Дж}$. На сколько увеличится средняя кинетическая энергия молекул этого газа после увеличения средней квадратичной скорости его молекул в $2$ раза? (В $10^{-22}\,\text{Дж}$; до целого числа.)
}{
\smallskip\noindent\textbf{Решение}\par
Средняя кинетическая энергия и среднеквадратичная скорость молекул связаны соотношением $E= \frac{m_0 v ^{2}}{2} .$ При увеличении средней квадратичной скорости в $2$ раза кинетическая энергия увеличится в $4$ раза, то есть станет равной $4\cdot 56$ · $10^{-22}=224\cdot 10^{-22}\,\text{Дж}$. Таким образом, кинетическая энергия увеличилась на $224\cdot 10^{-22}$ - $56\cdot 10^{-22}=168$ · $10^{-22}\,\text{Дж}$.

Ответ: $168$.
}

% Source sdamgia-20018: https://phys-ege.sdamgia.ru/problem?id=20018
\BankTask{sdamgia-20018}{
При уменьшении абсолютной температуры на $600\,\text{К}$ средняя кинетическая энергия теплового движения молекул аргона уменьшилась в $4$ раза. Какова конечная температура аргона? (В кельвинах.)
}{
\smallskip\noindent\textbf{Решение}\par
Средняя кинетическая энергия зависит от температуры в соответствии с формулой $\overline{E}_{\text{к}}= \frac{3}{2} kT.$ По условию кинетическая энергия уменьшилась в $4$ раза. Следовательно, температура также уменьшилась в $4$ раза, то есть $T_1=4T_2.$ По условию задачи температура уменьшилась на $\Delta T=T_1 - T_2.$ Тогда $\Delta T=4T_2 - T_2=3T_2.$ Откуда $T_2= \frac{\Delta T}{3} = \frac{600}{3} =200 \text{К}.$

Ответ: $200$.
}

% Source sdamgia-50742: https://phys-ege.sdamgia.ru/problem?id=50742
\BankTask{sdamgia-50742}{
В двух сосудах одинакового объёма находятся разреженные газы. В первом сосуде находится $2\,\text{моль}$ неона при температуре $127^\circ$C, во втором сосуде находится $2\,\text{моль}$ аргона при температуре $200$ $K$.

$1$. Абсолютная температура газа во втором сосуде выше, чем в первом.

$2$. Концентрация газа во втором сосуде меньше, чем в первом.

$3$. Отношение средней кинетической энергии теплового движения молекул неона к средней кинетической энергии теплового движения молекул аргона равно $2$.

$4$. Давление аргона больше неона.

$5$. Среднеквадратичная скорость молекул газа в первом сосуде больше, чем во втором.
}{
\smallskip\noindent\textbf{Решение}\par
$1$) Неверно. Абсолютная температура газа в первом сосуде равна $127^\circ$C + $273=400\,\text{К}$.

$2$) Неверно. Концентрация газа $n= \frac{N}{V} = \frac{\nu N_A}{V} .$ При равных объемах и количестве вещества концентрация одинакова.

$3$) Верно. Отношение $\frac{\overline{E}_{\text{н}}}{\overline{E}_{\text{а}}} = \frac{T_{\text{н}}}{T_{\text{а}}} = \frac{400}{200} =2.$

$4$) Неверно. Давление газа $p=nkT.$ При равных концентрациях температура неона больше, следовательно, давление неона больше.

$5$) Верно. Средняя квадратичная скорость молекул $\overline v = \sqrt{\frac{3RT}{M}} .$ Тогда

$\frac{\overline{v}_{\text{н}}}{\overline{v}_{\text{а}}}=\sqrt{\frac{T_{\text{н}}M_{\text{а}}}{T_{\text{а}}M_{\text{н}}}}=\sqrt{\frac{400\cdot40\cdot10^{-3}}{200\cdot20\cdot10^{-3}}}=2>1.$

Ответ: $35$.
}

% Source sdamgia-23234: https://phys-ege.sdamgia.ru/problem?id=23234
\BankTask{sdamgia-23234}{
В сосуде с жесткими стенками находится в равновесном состоянии смесь одного моля гелия и одного моля аргона. Температуру смеси повысили.

$1$. Среднеквадратичные скорости молекул гелия и молекул аргона в равновесном состоянии одинаковы.

$2$. Средние кинетические энергии поступательного теплового движения молекул гелия и молекул аргона в равновесном состоянии одинаковы.

$3$. В результате повышения температуры внутренняя энергия гелия увеличилась больше, чем внутренняя энергия аргона.

$4$. После повышения температуры парциальное давление гелия в сосуде увеличилось.

$5$. После повышения температуры давление в сосуде увеличилось.
}{
\smallskip\noindent\textbf{Решение}\par
$1$. Неверно. Средняя квадратичная скорость идеального газа равна $\overline v = \sqrt{\frac{3RT}{M}} .$ Поскольку молярная масса аргона больше, чем молярная масса гелия, то средняя квадратичная скорость аргона меньше средней квадратичной скорости гелия.

$2$. Верно. Средняя кинетическая энергия идеального газа определяется формулой $\overline{E}_k= \frac{3}{2} kT.$ При установлении равновесной температуры средние кинетические энергии газов одинаковы.

$3$. Неверно. Изменение внутренней энергии идеального газа равно $\Delta U = \frac{3}{2} \nu R\Delta T.$ По условию количество вещества у газов одинаковые, оба газа нагрелись на одну и ту же температуру. Следовательно, изменение внутренней энергии обоих газов одинаковое.

$4$. Верно. Парциальное давление гелия в сосуде увеличилось, так как увеличилась его температура при неизменном объеме. 

$5$. Верно. Давление газа в сосуде равно сумме парциальных давлений аргона и гелия. При повышении температуры парциальные давления обоих газов увеличатся, а значит, увеличится и давление газа в сосуде.

Ответ: $245$.
}

% Source sdamgia-3695: https://phys-ege.sdamgia.ru/problem?id=3695
\BankTask{sdamgia-3695}{
Среднеквадратичная скорость молекул идеального одноатомного газа, заполняющего закрытый сосуд, равна $\overline v =450 \text{м}/\text{с}.$ Как и на сколько изменится среднеквадратичная скорость молекул этого газа, если давление в сосуде вследствие охлаждения газа уменьшить на $19$\%?
}{
\smallskip\noindent\textbf{Решение}\par
Среднеквадратичная скорость молекул идеального газа с молярной массой $M$ при температуре $T$ равна $\overline v = \sqrt{\frac{3RT}{M}} .$

Согласно уравнению Клапейрона — Менделеева:

$pV= \frac{m}{M} RT,$

где $\text{р}$ — давление газа, $V$ — объем сосуда, $m$ — масса газа. Из этих выражений следует, что $\overline v = \sqrt{\frac{3pV}{m}} .$ Тогда начальная и конечная среднеквадратичная скорости равны $\overline v _1= \sqrt{\frac{3p _1V}{m}}$ и $\overline v _2= \sqrt{\frac{3p _2V}{m}} ,$ здесь учтено, что изменение давления в сосуде происходит при неизменном объеме (сосуд закрытый).

Согласно условию задачи, $p_2=p_1 - 0{,}19p_1=0{,}81p_1.$ Следовательно:

$\overline v _2= \sqrt{\frac{3p _2V}{m}} = \sqrt{\frac{3 \cdot 0{,}81p_1V}{m}} =\overline v _1 \sqrt{0{,}81}=0{,}9\overline v _1.$

Отсюда следует, что изменение среднеквадратичной скорости молекул:

$\Delta \overline v =\overline v _2 - \overline v _1=0{,}9\overline v _1 - \overline v _1= - 0{,}1\overline v _1= - 0{,}1 \cdot 450\text{м}/\text{с}= - 45 \text{м}/\text{с}.$

Таким образом, среднеквадратичная скорость молекул газа уменьшится на $45\,\text{м/с}$.

Ответ: среднеквадратичная скорость молекул газа уменьшится на $45\,\text{м/с}$.

Приведем другое решение.

Запишем основное уравнение МКТ, для первого и второго состояний газа:

$p_1= \frac{1}{3} m_0n ( \overline v ^{2} ) _1,p_2= \frac{1}{3} m_0n ( \overline v ^{2} ) _2.$

Объем сосуда и число молекул в нем не изменяются, следовательно, концентрация остается неизменной. Получаем:

$\frac{p_1}{p_2} = \frac{( \overline v ^{2} ) _1}{( \overline v ^{2} ) _2} ,$ откуда $\overline v _2=\overline v _1 \sqrt{\frac{p_2}{p_1}} ,$

$\overline v _2=450 \sqrt{\frac{0{,}81p_1}{p_1}} =405\text{м}/\text{с}.$

Откуда $\overline v _2 - \overline v _1= - 45\text{м}/\text{с}.$

Ответ: $- 45\text{м}/\text{с}.$
}

% Source sdamgia-19683: https://phys-ege.sdamgia.ru/problem?id=19683
\BankTask{sdamgia-19683}{
В некотором процессе $1$-$2$, который проводился с одним молем идеального одноатомного газа, среднеквадратичная скорость молекул газа возросла от $u_{1}=350\,\text{м/с}$ до $u_{2}=380\,\text{м/с}$. При этом давление $p$ газа в процессе оставалось постоянным, и газ совершил работу $A=292\,\text{Дж}$. Чему равна молярная масса $\mu$ этого газа?
}{
\smallskip\noindent\textbf{Решение}\par
$1$. Согласно молекулярно-кинетической теории идеального газа, среднеквадратичная скорость его молекул связана с температурой газа формулой:

$u = \sqrt{\frac{3RT}{\mu}} ,$ откуда $T = \frac{\mu u ^{2}}{3R} .$

$2$. Работа газа при изобарическом (согласно условию) процессе равна $A = p ( V_2 - V_1 ) .$

$3$. Из уравнения Клапейрона — Менделеева для одного моля идеального газа получаем, что $p ( V_2 - V_1 ) = R ( T_2 - T_1 ) = A.$

$4$. Таким образом, из записанных уравнений следует, что $A = \frac{\mu}{3} \cdot ( u ^{2} _2 - u ^{2} _1 ) ,$ откуда:

$\mu = \frac{3A}{( u ^{2} _2 - u ^{2} _1 )} = \frac{3 \cdot 292}{30 \cdot 730} = 0{,}040\text{кг}/\text{моль} = 40\text{г}/\text{моль}.$

Ответ: $40\,\text{г}$/моль.
}

% Source sdamgia-29808: https://phys-ege.sdamgia.ru/problem?id=29808
\BankTask{sdamgia-29808}{
С одним молем неона провели процесс, в котором газ из начального состояния с давлением $p_{1}=10^{5}\,\text{Па}$ и объемом $V_{1}=25\,\text{л}$ перешел в конечное состояние с объемом $V_{2}=50\,\text{л}$. Давление при этом изменялось по линейному закону $p=p_0 + \alpha V$ где $p_{0}$ — некоторая постоянная величина, $\alpha = \frac{p_1}{2V_1} .$ На какую величину изменилась при этом среднеквадратичная скорость движения атомов неона?
}{
\smallskip\noindent\textbf{Решение}\par
$1$. Согласно уравнению состояния Менделеева — Клапейрона для $\nu=1\text{моля}$ идеального одноатомного газа при температуре $T$ имеем: $pV= v RT.$

$2$. Среднеквадратичная скорость атомов массой $m$ идеального одноатомного газа равна:

$\overline v = \sqrt{\frac{3kT}{m}} = \sqrt{\frac{3RT}{\mu}} = \sqrt{\frac{3pV}{\mu \nu}} .$

$3$. Отсюда в начальном состоянии имеем: 

$\overline{v}_1=\sqrt{\frac{3p_1V_1}{\mu\nu}}=\sqrt{\frac{3\cdot10^5\cdot25\cdot10^{-3}}{20\cdot10^{-3}}}\approx612{,}4\text{м/с}.$

$4$. В конце процесса:

$p_2=p_1 + \Delta p=p_1 + \alpha \Delta V=p_1 + \frac{p_1}{2V_1} ( V_2 - V_1 ) = \frac{p_1 ( V_1 + V_2 )}{2V_1} =1{,}5 \cdot 10 ^{5} \text{Па}.$

$5$. Таким образом:

$\overline{v}_2=\sqrt{\frac{3p_2V_2}{\mu\nu}}=\sqrt{\frac{3\cdot1{,}5\cdot10^5\cdot50\cdot10^{-3}}{20\cdot10^{-3}}}\approx1060{,}7\text{м/с}.$

$6$. Окончательно получаем: 

$\overline{v}_2-\overline{v}_1=\sqrt{\frac{3p_2V_2}{\mu\nu}}-\sqrt{\frac{3p_1V_1}{\mu\nu}}\approx1060{,}7-612{,}4\approx448\text{м/с}.$

Ответ: $\overline v _2 - \overline v _1\approx448\text{м}/\text{с}.$
}


% Source 403D76: https://ege.fipi.ru/bank/questions.php?search=1&proj=BA1F39653304A5B041B656915DC36B38&qid=403D76
\BankTask{403D76}{
В двух сосудах одинакового объёма находятся разреженные газы.
В первом сосуде находится \(2\,\text{моль}\) гелия при температуре \(127\,^{\circ}\mathrm{C}\), во втором сосуде находится \(1\,\text{моль}\) аргона при температуре \(300\,\text{К}\).
\begin{enumerate}
\item Абсолютная температура газа во втором сосуде выше, чем в первом.
\item Среднеквадратичная скорость молекул газа в первом сосуде больше, чем во втором.
\item Давления газов в сосудах одинаковы.
\item Концентрация газа в первом сосуде в \(2\) раза меньше, чем во втором.
\item Отношение средней кинетической энергии теплового движения молекул аргона к средней кинетической энергии теплового движения молекул гелия равно \(0{,}75\).
\end{enumerate}
}{
\smallskip\noindent\textbf{Решение от GPT}\par
\(T_1=127+273=400\,\text{К}\), \(T_2=300\,\text{К}\).
\begin{enumerate}
\item Неверно: \(T_2<T_1\).
\item Верно: \(v_{\text{кв}}=\sqrt{3RT/M}\); у гелия температура выше, а молярная масса меньше, поэтому \(v_{\text{кв},1}>v_{\text{кв},2}\).
\item Неверно: \(p_1/p_2=\nu_1T_1/(\nu_2T_2)=8/3\).
\item Неверно: при одинаковых объёмах \(n_1/n_2=\nu_1/\nu_2=2\), то есть концентрация в первом сосуде больше.
\item Верно: \(\overline E=3kT/2\), поэтому \(\overline E_{\mathrm{Ar}}/\overline E_{\mathrm{He}}=300/400=0{,}75\).
\end{enumerate}
Ответ: \(25\).
}

% Source 6907: https://phys-ege.sdamgia.ru/problem?id=6907
\BankTask{6907}{
Цилиндрический сосуд разделен неподвижной теплоизолирующей перегородкой. В одной части сосуда находится кислород, в другой  — водород, концентрации газов одинаковы. Давление кислорода в $2$ раза больше давления водорода. Чему равно отношение средней кинетической энергии теплового движения молекул кислорода к средней кинетической энергии теплового движения молекул водорода?
}{
\smallskip\noindent\textbf{Решение}\par
Запишем соотношение между давлением и средней кинетической энергий молекул: $p= \frac{2}{3} n\overline E_{\,\text{к}},$ где n  — концентрация газа.
При условии равенства концентраций кислорода и водорода получим отношение средних кинетических энергий: 

\[\frac{\overline E_{\,\text{к},\,\text{к}}}{\overline E_{\,\text{к},\,\text{в}}} = \frac{p_{\,\text{к}}}{p_{\,\text{в}}} =2.\]

Ответ: $2$.
}

\end{document}
